%% Anexo A — Infraestructura de trabajo y estructura del repositorio
%% Fragmento: no se compila solo. Lo incluye docs/entrega/entrega-fase1.tex.

\section{Anexo A. Infraestructura de trabajo}

La constitución de una consultora de software no se agota en su identidad y sus
roles: exige un lugar donde el trabajo quede registrado y sea verificable. Este
anexo describe la infraestructura con la que Legasoft inicia operaciones, la
organización del repositorio que la sostiene y las convenciones de gestión de
configuración que el equipo se compromete a respetar durante el \lgSemestre.

\subsection{Repositorio institucional}

Todo el trabajo de la consultora se concentra en un repositorio Git único,
público y ya operativo:

\begin{nota}[Dirección del repositorio]
\href{\lgRepo}{\lgRepoCorto}
\end{nota}

La decisión de mantener un solo repositorio —en lugar de separar documentación y
código— responde al principio de trazabilidad declarado en la constitución: un
requisito, la decisión de arquitectura que lo resuelve, el código que lo
implementa y la prueba que lo verifica quedan en la misma historia de cambios y
pueden recorrerse en ambos sentidos.

\subsection{Estructura de carpetas}

La estructura refleja la organización de la empresa: cada carpeta de primer nivel
corresponde a un ámbito de responsabilidad de los definidos en la sección de
roles, y cada una incluye su propio \ruta{README.md} con el detalle de lo que
contiene.

\begin{center}
\begin{tabularx}{\linewidth}{@{}L{3.5cm} Y L{3.3cm}@{}}
\toprule
\sffamily\bfseries\color{lgPrimario}Carpeta &
\sffamily\bfseries\color{lgPrimario}Contenido &
\sffamily\bfseries\color{lgPrimario}Responsable \\
\midrule
\ruta{docs/institucional/} & Constitución de la empresa, acta de constitución y documentación institucional. & Dirección de Proyecto \\
\ruta{docs/entrevistas/}   & Minutas y transcripciones de entrevistas de relevamiento (ENT / VAL). & \lgDirectorRol \\
\ruta{docs/srs/}           & Especificación de requisitos (SRS), historias de usuario y criterios de aceptación. & \lgDirectorRol \\
\ruta{docs/trazabilidad/}  & Matrices de trazabilidad entre necesidades, requisitos, arquitectura y pruebas. & \lgDirectorRol \\
\ruta{docs/minutas/}       & Minutas de reunión, acuerdos con el cliente y evidencia de validación. & \lgDirectorRol \\
\ruta{docs/arquitectura/}  & Documento de arquitectura y decisiones registradas (ADR). & \lgArquitectoRol \\
\ruta{docs/calidad/}       & Estrategia de calidad, plan de pruebas y casos de prueba. & \lgQARol \\
\ruta{docs/latex/}         & Estilo compartido de todos los documentos (\ruta{.sty}). & Compartida \\
\ruta{docs/cv/}            & Hojas de vida de los integrantes del equipo. & Dirección de Proyecto \\
\ruta{docs/entrega/}       & Documentos de entrega consolidados que se presentan al docente. & Dirección de Proyecto \\
\ruta{modelos/c4/}         & Diagramas del modelo C4: contexto, contenedores y componentes. & \lgArquitectoRol \\
\ruta{modelos/uml/}        & Diagramas UML: casos de uso, clases y secuencia. & \lgArquitectoRol \\
\ruta{modelos/bpmn/}       & Modelos de procesos de negocio AS-IS y TO-BE. & \lgDirectorRol \\
\ruta{src/backend/}        & Código del servidor, APIs y lógica de negocio. & \lgArquitectoRol \\
\ruta{src/frontend/}       & Código de la interfaz de usuario. & \lgQARol \\
\ruta{tests/unit/}         & Pruebas unitarias. & Compartida \\
\ruta{tests/integration/}  & Pruebas de integración. & \lgQARol \\
\ruta{tests/acceptance/}   & Pruebas de aceptación contra los criterios acordados. & \lgQARol \\
\ruta{.github/workflows/}  & Flujos de integración y entrega continua. & \lgQARol \\
\bottomrule
\end{tabularx}
\captionof{table}{Estructura del repositorio y responsable de cada carpeta.}
\end{center}

\subsection{Gestión de la documentación}

La documentación se escribe en \LaTeX{} y se compila a PDF con un único comando
(\ruta{make}). Esta decisión persigue tres efectos concretos:

\begin{itemize}
  \item \textbf{Formato consistente.} Todos los documentos comparten el estilo
        \ruta{docs/latex/legasoft.sty}, de modo que la portada, la numeración, los
        encabezados y las tablas son idénticos en cualquier entregable, sin
        depender de que cada integrante recuerde aplicarlos.
  \item \textbf{Control de versiones real.} El texto fuente es plano, por lo que
        Git muestra qué frase cambió, quién la cambió y cuándo, y permite revisar
        un documento con el mismo mecanismo con que se revisa el código.
  \item \textbf{Identificación de entregables.} Cada documento declara su código,
        su versión y su estado, y abre con una tabla de control de versiones.
\end{itemize}

Los documentos aún no completados marcan en ámbar los datos que faltan
(\campo{así}), para que ninguna plantilla pueda confundirse con un documento
aprobado.

\subsection{Nomenclatura de los documentos}

Cada entregable recibe un código estable que lo identifica con independencia de su
nombre de archivo:

\begin{center}
\begin{tabularx}{\linewidth}{@{}L{3.2cm} Y@{}}
\toprule
\sffamily\bfseries\color{lgPrimario}Código &
\sffamily\bfseries\color{lgPrimario}Documento \\
\midrule
\texttt{LGS-ENT-NNN} & Documento de entrega consolidado. \\
\texttt{LGS-INST-NNN} & Documentación institucional de la consultora. \\
\texttt{LGS-REQ-NNN} & Especificación de requisitos y documentos asociados. \\
\texttt{LGS-ARQ-NNN} & Documento de arquitectura. \\
\texttt{LGS-ADR-NNN} & Decisión de arquitectura registrada. \\
\texttt{LGS-CAL-NNN} & Estrategia de calidad y plan de pruebas. \\
\texttt{LGS-MIN-NNN} & Minuta de reunión. \\
\texttt{LGS-PRY-NNN} & Documentos propios de un proyecto con cliente. \\
\texttt{LGS-RRHH-NNN} & Hojas de vida y gestión de recursos humanos. \\
\bottomrule
\end{tabularx}
\captionof{table}{Codificación de los documentos de Legasoft.}
\end{center}

\subsection{Control de versiones}

El equipo trabaja sobre la rama \ruta{main}, que debe permanecer siempre en un
estado presentable. El trabajo en curso se realiza en ramas temáticas con el
prefijo del ámbito al que pertenece —\ruta{docs/}, \ruta{feat/}, \ruta{fix/} o
\ruta{test/}— y se integra mediante \emph{pull request} revisada por un
integrante distinto del autor. Los mensajes de \emph{commit} siguen la convención
de \emph{Conventional Commits} (\ruta{tipo: descripción en imperativo}), de modo
que la historia del repositorio se lea como un registro de decisiones y no como
una sucesión de guardados.

Los PDF son artefactos de compilación y no se versionan, con una excepción
deliberada: el documento de entrega consolidado sí se conserva en el repositorio,
porque constituye la evidencia de lo que se presentó y en qué fecha.

\subsection{Herramientas de trabajo}

\begin{center}
\begin{tabularx}{\linewidth}{@{}L{3.6cm} Y@{}}
\toprule
\sffamily\bfseries\color{lgPrimario}Herramienta &
\sffamily\bfseries\color{lgPrimario}Uso en la consultora \\
\midrule
Git y GitHub & Control de versiones, revisión entre pares y alojamiento del repositorio institucional. \\
Jira & Cronograma, tablero de tareas, seguimiento del avance y gestión del \emph{backlog}. \\
\LaTeX{} con Tectonic & Redacción y compilación de toda la documentación a PDF. \\
GitHub Actions & Integración y entrega continua: compilación de la documentación y ejecución de pruebas. \\
Reuniones registradas & Cada reunión con el cliente o interna deja minuta en \ruta{docs/minutas/}. \\
\bottomrule
\end{tabularx}
\captionof{table}{Herramientas que sostienen el proceso de trabajo.}
\end{center}

\begin{nota}[Alcance de esta infraestructura]
Las carpetas \ruta{src/} y \ruta{tests/} están creadas y documentadas, pero
todavía vacías de código: las tecnologías de frontend, backend y pruebas se
decidirán una vez conocido el cliente y quedarán registradas como decisiones de
arquitectura (ADR) en \ruta{docs/arquitectura/}. Declararlas ahora contradiría el
principio de intervención de la consultora, según el cual cada proyecto debe
justificar la tecnología que incorpora.
\end{nota}
